% Options for packages loaded elsewhere
\PassOptionsToPackage{unicode}{hyperref}
\PassOptionsToPackage{hyphens}{url}
\documentclass[
]{article}
\usepackage{xcolor}
\usepackage[margin=1in]{geometry}
\usepackage{amsmath,amssymb}
\setcounter{secnumdepth}{5}
\usepackage{iftex}
\ifPDFTeX
  \usepackage[T1]{fontenc}
  \usepackage[utf8]{inputenc}
  \usepackage{textcomp} % provide euro and other symbols
\else % if luatex or xetex
  \usepackage{unicode-math} % this also loads fontspec
  \defaultfontfeatures{Scale=MatchLowercase}
  \defaultfontfeatures[\rmfamily]{Ligatures=TeX,Scale=1}
\fi
\usepackage{lmodern}
\ifPDFTeX\else
  % xetex/luatex font selection
\fi
% Use upquote if available, for straight quotes in verbatim environments
\IfFileExists{upquote.sty}{\usepackage{upquote}}{}
\IfFileExists{microtype.sty}{% use microtype if available
  \usepackage[]{microtype}
  \UseMicrotypeSet[protrusion]{basicmath} % disable protrusion for tt fonts
}{}
\makeatletter
\@ifundefined{KOMAClassName}{% if non-KOMA class
  \IfFileExists{parskip.sty}{%
    \usepackage{parskip}
  }{% else
    \setlength{\parindent}{0pt}
    \setlength{\parskip}{6pt plus 2pt minus 1pt}}
}{% if KOMA class
  \KOMAoptions{parskip=half}}
\makeatother
\usepackage{longtable,booktabs,array}
\usepackage{calc} % for calculating minipage widths
% Correct order of tables after \paragraph or \subparagraph
\usepackage{etoolbox}
\makeatletter
\patchcmd\longtable{\par}{\if@noskipsec\mbox{}\fi\par}{}{}
\makeatother
% Allow footnotes in longtable head/foot
\IfFileExists{footnotehyper.sty}{\usepackage{footnotehyper}}{\usepackage{footnote}}
\makesavenoteenv{longtable}
\usepackage{graphicx}
\makeatletter
\newsavebox\pandoc@box
\newcommand*\pandocbounded[1]{% scales image to fit in text height/width
  \sbox\pandoc@box{#1}%
  \Gscale@div\@tempa{\textheight}{\dimexpr\ht\pandoc@box+\dp\pandoc@box\relax}%
  \Gscale@div\@tempb{\linewidth}{\wd\pandoc@box}%
  \ifdim\@tempb\p@<\@tempa\p@\let\@tempa\@tempb\fi% select the smaller of both
  \ifdim\@tempa\p@<\p@\scalebox{\@tempa}{\usebox\pandoc@box}%
  \else\usebox{\pandoc@box}%
  \fi%
}
% Set default figure placement to htbp
\def\fps@figure{htbp}
\makeatother
\setlength{\emergencystretch}{3em} % prevent overfull lines
\providecommand{\tightlist}{%
  \setlength{\itemsep}{0pt}\setlength{\parskip}{0pt}}
\usepackage{booktabs}
\usepackage{caption}
\usepackage{float}
\renewcommand{\familydefault}{\sfdefault}
\usepackage{placeins}
\usepackage{caption}
\captionsetup[table]{position=bottom}
\usepackage{bookmark}
\IfFileExists{xurl.sty}{\usepackage{xurl}}{} % add URL line breaks if available
\urlstyle{same}
\hypersetup{
  pdftitle={Linear Independence of Variables in INTOXICATE-ER},
  pdfauthor={Yash Jaggi},
  hidelinks,
  pdfcreator={LaTeX via pandoc}}

\title{Linear Independence of Variables in INTOXICATE-ER}
\author{Yash Jaggi}
\date{}

\begin{document}
\maketitle

\section*{Realtionship Between Gender and Xenobiotic}\label{realtionship-between-gender-and-xenobiotic}
\addcontentsline{toc}{section}{Realtionship Between Gender and Xenobiotic}

\begin{table}[!h]
\centering
\caption{\label{tab:table1}\textbf{No association between gender and type of xenobiotic.}. Cells denote case count (row percentage). F, biological female. M, male. See X for description of exposure categories.}
\centering
\begin{tabular}[t]{lcccccc}
\toprule
\multicolumn{1}{c}{ } & \multicolumn{2}{c}{Biological Sex} & \multicolumn{4}{c}{ } \\
\cmidrule(l{3pt}r{3pt}){2-3}
\textbf{} & \textbf{\textit{F}} & \textbf{\textit{M}} & \textbf{Total} & \textbf{p-value} & \textbf{\textbf{Row vs Rest p}} & \textbf{\textbf{Row vs Rest p (adj)}}\\
\midrule
Exposure Category &  &  &  &  &  & \\
\hspace{1em}\em{Combination} & 20 (41\%) & 29 (59\%) & 49 (100\%) &  & 0.266 & 0.531\\
\hspace{1em}\em{Unknown} & 13 (50\%) & 13 (50\%) & 26 (100\%) &  & 1.000 & 1.000\\
\hspace{1em}\em{Analgesic} & 16 (64\%) & 9 (36\%) & 25 (100\%) &  & 0.153 & 0.407\\
\hspace{1em}\em{Antidepressants} & 17 (74\%) & 6 (26\%) & 23 (100\%) &  & 0.018 & 0.146\\
\addlinespace
\hspace{1em}\em{Street Drugs} & 9 (39\%) & 14 (61\%) & 23 (100\%) &  & 0.450 & 0.545\\
\hspace{1em}\em{CO, As, CN} & 7 (32\%) & 15 (68\%) & 22 (100\%) &  & 0.145 & 0.407\\
\hspace{1em}\em{Sedatives} & 7 (64\%) & 4 (36\%) & 11 (100\%) &  & 0.476 & 0.545\\
\hspace{1em}\em{Alcohol} & 3 (30\%) & 7 (70\%) & 10 (100\%) &  & 0.374 & 0.545\\
Total & 92 (49\%) & 97 (51\%) & 189 (100\%) &  &  & \\
\bottomrule
\multicolumn{7}{l}{\rule{0pt}{1em}\textsuperscript{1} NA}\\
\end{tabular}
\end{table}

\begin{verbatim}
## 
## Call:
## glm(formula = `Dichotomous Outcome` ~ `INTOXICATE SCORE` + Gender, 
##     family = binomial, data = train_data)
## 
## Coefficients:
##                    Estimate Std. Error z value Pr(>|z|)    
## (Intercept)        -3.10696    0.66253  -4.690 2.74e-06 ***
## `INTOXICATE SCORE`  0.10510    0.04316   2.435   0.0149 *  
## GenderM             1.31201    0.61853   2.121   0.0339 *  
## ---
## Signif. codes:  0 '***' 0.001 '**' 0.01 '*' 0.05 '.' 0.1 ' ' 1
## 
## (Dispersion parameter for binomial family taken to be 1)
## 
##     Null deviance: 83.702  on 81  degrees of freedom
## Residual deviance: 70.138  on 79  degrees of freedom
##   (19 observations deleted due to missingness)
## AIC: 76.138
## 
## Number of Fisher Scoring iterations: 5
\end{verbatim}

\begin{verbatim}
##                          2.5 %     97.5 %
## (Intercept)        -4.56486943 -1.9306520
## `INTOXICATE SCORE`  0.02486071  0.1967328
## GenderM             0.13286220  2.6000521
\end{verbatim}

The sensitivity and specificity of the model on the training data are 0.706 and 0.508, respectively. On the test data, the sensitivity and specificity are 0.267 and 1, respectively.

\end{document}
